\chapter{Conclusion and Further Work}
\label{chap:conclusion}

\section{Summary}
\label{sec:summary}

The project asked whether PerMFL transfers to network intrusion detection, what
has to change for it to transfer, and whether the architectural claim that
motivates it holds on this kind of data.

It transfers, but not as published. The released implementation reproduces its
own headline figure to within 0.04 percentage points, and it does so while
carrying a defect that silently overrides a documented hyperparameter with an
unrelated loop count. That the published number is reached anyway is evidence
the defect was present when the results were produced.

The architectural claim does not hold in the region the method's own theory
permits. Removing the team tier entirely changes personalised macro F1 by
0.0002 on CICIDS2017 and by 0.0000 on the base paper's own EMNIST-10 setup,
against a measured noise floor of 0.0116. Assigning teams to match true attack
domains does not beat assigning them at random, and neither beats teams derived
from client updates during training.

The reason is arithmetic rather than empirical. A single parameter governs both
the device-to-team penalty and the weight the team model places on its members,
and those two roles want opposed values. At the setting that personalises well,
membership influences the team model by 3.3 per cent, so the team model is
effectively a copy of the global one. The method's own convergence condition
caps that influence below fifty per cent in any case.

Separating the two roles improves both models. Across three datasets and nine
heterogeneity settings the correction raises personalised macro F1 by 0.027 to
0.326 and global macro F1 by 0.113 to 0.705, significant at every point. It also
reduces across-seed variance in global accuracy by two to three orders of
magnitude, which holds even in the one configuration where the mean regressed.

One result runs against the project's expectations and is reported as such. The
benefit grows as clients become more similar, not less. The term the weight
controls is an average of team members, and an average of divergent members
carries little signal, so the correction mostly frees the model to use
information it already had.

\section{Against the objectives}
\label{sec:objectives-met}

The reproduction was achieved and the pipeline built for three datasets with
four partitioning schemes and a documented evaluation protocol. Heterogeneity
was swept from near-homogeneous to strongly skewed with every result anchored
against a floor and, where the protocol imposes one, a structural ceiling. The
architectural premise was tested directly, the mechanism behind the negative
result identified from the update rule, and a correction evaluated with paired
seeds and significance testing.

One objective is only partly met. The base implementation ships seven baseline
algorithms and the intention was to run them on the intrusion detection data.
Only the method under test carries the metric instrumentation this study
requires, and extending it to the others was not completed. The results
therefore position PerMFL against its own configurations and against a local
training baseline, not against the wider federated literature.

\section{Limitations}
\label{sec:limitations}

The absent baseline comparison is the principal limitation. A reader cannot tell
from this work whether the corrected method is competitive with FedAvg, pFedMe
or Ditto on the same data.

No comparison is made against the two most closely related methods, SCMoE-PFL
and CFMD-i, because neither publishes an implementation. This also means no
claim is made about relative computational cost, despite the substantial
practical difference in what the architectures hold: one model per tier against
per-cluster experts, a private model and a gating network per client.

Several exploratory results rest on single runs, and are labelled as such.

The heterogeneity measure captures divergence between label distributions only.
A partition with identical label distributions but different decision boundaries
would register as homogeneous.

Two of the three datasets are laboratory captures with scripted attacks, so
agreement between them is weaker evidence than agreement between independent
operational sources.

\section{Further work}
\label{sec:further-work}

\subsection{Completing the comparison}

Instrumenting the shipped baselines with the same metrics and running them on
the same partitions would place the corrected method in context. This is the
most valuable remaining work and it needs no new ideas, only the same
instrumentation applied to four more server classes.

\subsection{Tuning the decoupled weight}

The team weight's optimum is dataset dependent: 1.5 on CICIDS2017 against 3.0
on NSL-KDD. A principled rule relating it to the heterogeneity measure, rather
than a per-dataset sweep, would make the correction usable without tuning.

\subsection{Output-space clustering}

The derived team formation uses a parameter-space signal. CFMD-i additionally
uses an output-space divergence and switches between the two adaptively. Under
label skew two clients can hold near-identical weights and behave differently,
which the parameter-space signal misses. Whether the better signal changes the
conclusion is uncertain given that the term it feeds is weighted at 3.3 per
cent, but the argument is currently an argument rather than a measurement.

\subsection{Nested structure}

Every partition used here supplies structure at one level, so two of the three
tiers are redundant by construction. A nested partition, in which teams differ
by attack family and clients within a team differ by variant, would give each
tier something distinct to represent. It is the fairest available test of the
three-tier premise and no published work has run it.

\subsection{Quantity skew}

The per-client cap flattens client sizes into a narrow band. Real deployments
are unequal, and the base paper's own synthetic generator spans two orders of
magnitude. Removing the cap restores an axis of realism that this study
suppressed.

\subsection{Memory behaviour of the loader interface}

The loader interface materialises every client's features as Python lists,
roughly a gigabyte per process on CICIDS2017 before any training begins. This
constrains how many configurations can run concurrently and caused a
machine-level failure during this project. A tensor-backed interface would
remove the constraint, at the cost of changing a contract every loader in the
base implementation shares.

\section{Closing}
\label{sec:closing}

The most useful thing this project produces is not the correction. It is the
observation that a published method's central architectural claim can be tested
directly, cheaply, and negatively, and that the reason can be read off the
update rule once someone looks. The team tier contributes 0.0002 macro F1 in the
configuration its authors recommend, and the arithmetic explaining why fits in a
paragraph.

That the same analysis then yields a correction improving both models across
three datasets is a secondary result, though the more useful one to anyone
intending to deploy the method.
